\section{The ordinary fundamental-group question}\label{sec:ordinary}

The monodromy theorem concerns every complex linear image, but does
not by itself determine the ordinary group. The remaining assertion
can be stated without representation-theoretic terminology.

\begin{question}\label{q:full}
For every connected smooth special complex quasi-projective variety
$V$, does $\pi_1(V)$ have a finite-index subgroup $H$ satisfying
$[H,[H,H]]=1$?
\end{question}

This is \cite[Conjecture 2.12]{CDHP2026}. The group is the ordinary
topological fundamental group; boundary meridians have not been
quotiented out. The surfaces in Theorem~\ref{thm:example} satisfy the
proposed conclusion and hence do not give counterexamples to it.

The precise obstruction to an inference from Theorem~\ref{thm:completion}
is faithfulness. Put $\Gamma=\pi_1(V)$ and write $\Gamma^{\un}$ for
its rational unipotent completion. The canonical homomorphism
$\Gamma\to\Gamma^{\un}(\Q)$ can have a kernel. The theorem sends
$\gamma_3\Gamma$ to the identity without identifying its elements in
$\Gamma$. For example, a finitely generated group with finite
abelianization has trivial rational unipotent completion, since a
nonzero nilpotent Lie algebra has nonzero abelianization. This
shows that a nontrivial group can be entirely invisible to rational
unipotent completion. The specialness hypotheses used here supply no
injectivity statement for the completion map.

A known group-theoretic reduction shows what an additional geometric
input could accomplish. If the full group $\pi_1(V)$ were virtually
solvable, Rogov's theorem would make it virtually nilpotent
\cite[Theorem C, Corollary 6.6]{Rogov2025}. A finitely generated
virtually nilpotent group has a torsion-free nilpotent subgroup $H$ of
finite index, and $H$ admits a faithful rational unipotent
representation by Mal'cev's theorem
\cite[Chapters 1 and 5; Chapter 6, p.~100]{Segal1983}.
The finite \'etale cover corresponding to $H$ is again special. The CDY
quasi-Albanese inputs from Section~\ref{sec:linear} therefore apply to
that cover, and Theorem~\ref{thm:completion} gives $[H,[H,H]]=1$.
Virtual solvability of the full ordinary group is not established by
the representation bounds.

Theorem~\ref{thm:example} also prevents an induction that requires the
general quasi-Albanese fibre itself to be special. Its fundamental
group is nonabelian free, while its image in the total group is zero.
Controlling that actual image is a different geometric problem.
The completion theorem addresses its image in every canonical
unipotent quotient; control of the kernel invisible to these quotients
remains necessary for Question~\ref{q:full}.
